%
% teil6.tex -- Modifizierte Bessel-Funktionen
%
% (c) 2026 Gian Cavegn, OST Ostschweizer Fachhochschule
%
% !TEX root = ../../buch.tex
% !TEX encoding = UTF-8
%
\section{Modifizierte Bessel-Funktionen
\label{bessel:section:modifiziert}}
\kopfrechts{Modifizierte Bessel-Funktionen}
% Funktion dieses Abschnitts:
% Ergänzung zum bisher behandelten oszillatorischen Fall. Die
% modifizierte Bessel-DGL unterscheidet sich nur durch ein
% Vorzeichen, führt aber auf Funktionen mit exponentiellem statt
% oszillierendem Verhalten -- analog zum Verhältnis zwischen
% Kosinus und Cosinus Hyperbolicus.
%
% Empfohlene Position: nach teil3 (Eigenschaften), vor teil4
% (Lichtbeugung). In main.tex entsprechend einbinden.

In Abschnitt~\ref{bessel:section:herleitung} entstand die Bessel-DGL aus der Helmholtz-Gleichung mit reellem Wellenzahlquadrat $k^2 > 0$.
Bei stationären Diffusions- oder Wärmeleitungsproblemen sowie bei Wellenleitern im evaneszenten Bereich tritt jedoch effektiv ein \emph{negatives} Wellenzahlquadrat auf.
Das führt auf eine leicht abgeänderte Differentialgleichung, deren Lösungen sich grundlegend anders verhalten.

\subsection{Die modifizierte Bessel-Differentialgleichung
\label{bessel:subsection:mod-dgl}}
% Inhalt:
% - Ausgangspunkt: gleiche Helmholtz-Separation wie in teil1, aber
%   mit k^2 -> -k^2 (oder z-Abhängigkeit cosh/sinh statt cos/sin).
% - Substitution x = kr liefert die modifizierte Bessel-DGL.
% - Alternativ: x -> ix in der gewöhnlichen Bessel-DGL einsetzen.

Geht man bei der Separation der Helmholtz-Gleichung von einer $z$-Abhängigkeit der Form $\cosh(kz), \sinh(kz)$ statt $\cos(kz), \sin(kz)$ aus -- wie es etwa bei stationären Problemen oder bei Wellenleitern mit imaginärer Wellenzahl der Fall ist -- so wechselt das Vorzeichen des $k^2$-Terms in der Radialgleichung~\eqref{bessel:equation:radial}.
Nach der gleichen Substitution $x = kr$ erhält man die \emph{modifizierte Bessel'sche Differentialgleichung}
\begin{equation}
\boxed{\;
x^2 y''(x) + x y'(x) - (x^2 + \nu^2) y(x) = 0\;
}
\label{bessel:equation:mod-bessel-dgl}
\end{equation}
Im Vergleich mit~\eqref{bessel:equation:bessel-dgl} ist lediglich das Vorzeichen vor $x^2$ umgekehrt.
Formal lässt sich diese Gleichung auch durch die Substitution $x \to ix$ in der gewöhnlichen Bessel-DGL gewinnen.
% [TODO: kurzer Nachweis dieser Tatsache]

\subsection{Die modifizierte Bessel-Funktion erster Art
\label{bessel:subsection:Inu}}
% Inhalt:
% - Wieder Frobenius-Ansatz; die Rechnung läuft praktisch identisch
%   zu teil2 ab, nur dass das Vorzeichen in der Rekursion wechselt.
% - Endergebnis: alternierende Vorzeichen fallen weg.
% - Definition I_nu als Reihe; Beziehung I_nu(x) = i^{-nu} J_nu(ix).

Der Frobenius-Ansatz aus Abschnitt~\ref{bessel:section:loesung} lässt sich unverändert auf~\eqref{bessel:equation:mod-bessel-dgl} anwenden.
Die Indexgleichung $s^2 - \nu^2 = 0$ bleibt dieselbe; in der Rekursion~\eqref{bessel:equation:rekursion} wechselt jedoch das Vorzeichen, so dass die alternierenden Faktoren $(-1)^m$ entfallen.
Mit derselben Normierung wie in~\eqref{bessel:equation:Jnu} ergibt sich die \emph{modifizierte Bessel-Funktion erster Art}
\begin{equation}
\boxed{\;
I_\nu(x)
=
\sum_{m=0}^{\infty}
\frac{1}{m!\, \Gamma(\nu+m+1)}
\left(\frac{x}{2}\right)^{2m+\nu}
\;}
\label{bessel:equation:Inu}
\end{equation}
Sie steht mit der gewöhnlichen Bessel-Funktion in der Beziehung
\begin{equation}
I_\nu(x) = i^{-\nu}\, J_\nu(ix).
\label{bessel:equation:Inu-Jnu}
\end{equation}
Da in~\eqref{bessel:equation:Inu} alle Reihenglieder positiv sind, ist $I_\nu(x)$ für $x > 0$ streng monoton wachsend und hat keine Nullstellen.
Das Verhalten ist also exponentiell wachsend, nicht oszillatorisch.

\subsection{Die modifizierte Bessel-Funktion zweiter Art
\label{bessel:subsection:Knu}}
% Inhalt:
% - Wie bei Y_nu: bei ganzzahligem nu sind I_nu und I_{-nu} linear
%   abhängig. Eine zweite, unabhängige Lösung wird gebraucht.
% - Konvention: die Macdonald-Funktion K_nu.
%   K_nu(x) = (pi/2) [I_{-nu}(x) - I_nu(x)] / sin(nu pi),
%   bei ganzzahligem nu als Grenzwert.
% - Wichtig: K_nu(x) -> 0 für x -> infinity, d.h. sie ist die
%   im Unendlichen abklingende Lösung.
% - Singularität bei x = 0.

Die zweite linear unabhängige Lösung der modifizierten Bessel-DGL wird als \emph{Macdonald-Funktion} oder \emph{modifizierte Bessel-Funktion zweiter Art} bezeichnet und ist durch
\begin{equation}
K_\nu(x)
=
\frac{\pi}{2}\,
\frac{I_{-\nu}(x) - I_\nu(x)}{\sin(\nu\pi)}
\label{bessel:equation:Knu}
\end{equation}
definiert, wobei für ganzzahliges $\nu$ der entsprechende Grenzwert zu bilden ist.
% [TODO: Bemerkung zum Grenzwert bei ganzzahligem nu, analog zu Y_nu]
Im Gegensatz zu $I_\nu$ klingt $K_\nu(x)$ für $x \to \infty$ exponentiell ab, ist dafür aber im Ursprung singulär.
In physikalischen Anwendungen wählt man $I_\nu$ als reguläre Lösung im Inneren eines Zylinders und $K_\nu$ als abklingende Lösung im Aussenraum.

\subsection{Asymptotisches Verhalten
\label{bessel:subsection:mod-asymptotik}}
% Inhalt:
% - I_nu(x) ~ (x/2)^nu / Gamma(nu+1) für x -> 0
% - I_nu(x) ~ e^x / sqrt(2 pi x) für x -> infty
% - K_nu(x) ~ sqrt(pi/(2x)) e^{-x} für x -> infty
% - K_0(x) ~ -ln(x) für x -> 0 (logarithmische Singularität)
% - Sonst K_nu(x) ~ (Gamma(nu)/2)(2/x)^nu für nu > 0.

Aus der Reihendarstellung~\eqref{bessel:equation:Inu} folgt für kleine Argumente
\begin{equation}
I_\nu(x)
\sim
\frac{1}{\Gamma(\nu+1)}\left(\frac{x}{2}\right)^\nu
\qquad (x \to 0).
\label{bessel:equation:Inu-klein}
\end{equation}
Für grosse Argumente gelten die asymptotischen Formeln
\begin{align}
I_\nu(x) &\sim \frac{e^x}{\sqrt{2\pi x}}
\qquad (x \to \infty),
\label{bessel:equation:Inu-gross} \\
K_\nu(x) &\sim \sqrt{\frac{\pi}{2x}}\, e^{-x}
\qquad (x \to \infty).
\label{bessel:equation:Knu-gross}
\end{align}
Während $J_\nu$ und $Y_\nu$ also wie gedämpfte Schwingungen wirken (vgl.~\eqref{bessel:equation:asymptotik-gross}), verhalten sich $I_\nu$ und $K_\nu$ wie exponentiell wachsende beziehungsweise abklingende Funktionen.
Das Begriffspaar gewöhnliche/modifizierte Bessel-Funktionen ist damit vollständig analog zum Paar trigonometrische/hyperbolische Funktionen.
% [TODO: Bild oder Tabelle, die J_0, Y_0 vs. I_0, K_0 visualisiert]

\subsection{Anwendungsbeispiel: Stationäre Wärmeleitung im Zylinder
\label{bessel:subsection:mod-anwendung}}
% Optionaler Anwendungsabschnitt -- kann auch weggelassen werden,
% falls das Kapitel sonst zu lang wird.
%
% Inhalt:
% - Stationäre Wärmeleitungsgleichung: Delta T = 0 in einem
%   unendlich langen Zylinder, mit gegebenen Randwerten an
%   Mantel (z.B. T(R, z) = f(z) periodisch in z).
% - Separationsansatz T = R(r) Z(z) mit Z = exp(ikz).
% - Radialgleichung wird zur modifizierten Bessel-DGL.
% - Im Inneren (regulär bei r=0) wählt man I_0(kr); im Aussenraum
%   K_0(kr).
% - Damit lassen sich konkrete Randwertprobleme lösen.

Als kurzes Anwendungsbeispiel betrachten wir die stationäre Temperaturverteilung $T(r, z)$ in einem unendlich langen, rotationssymmetrischen Zylinder.
% [TODO: Skizze des Aufbaus mit Zylinder, Koordinatensystem, Randbedingungen]
Die stationäre Wärmeleitungsgleichung lautet
\begin{equation}
\Delta T = 0,
\end{equation}
und mit dem Separationsansatz $T(r, z) = R(r) e^{ikz}$ ergibt sich für den Radialteil genau die modifizierte Bessel-DGL.
Die im Inneren ($r < R_0$) reguläre Lösung ist $I_0(kr)$, im Aussenraum ($r > R_0$) wählt man die abklingende Lösung $K_0(kr)$.
Damit lassen sich Randwertprobleme der stationären Wärmeleitung in zylindrischer Geometrie explizit lösen.
% [TODO: vollständige Durchrechnung eines konkreten Beispiels]

